\section{Setting, hypotheses and exact identities}\label{sec:setting}

\subsection{Runs and the occupation count}\label{sec:runs}

The walk takes $N$ steps and each step has a sign, $\plus$ or $\minus$. The signs
come in maximal runs of equal sign, and consecutive runs have opposite signs. All
runs after the first are independent. A run of sign $\sgn\in\{\plus,\minus\}$ has a
length distributed as $X_\sgn$, a random variable on $\{1,2,\dots\}$ with mass
function $f_\sgn(n)=\Prob(X_\sgn=n)$ and tail
\[
 T_\sgn(x)=\Prob(X_\sgn\ge x),\qquad x\in\R .
\]
We write $\osgn$ for the sign opposite to $\sgn$. The first run is the initial run.
It has a sign $\Sigma$ and a length $D$, the pair $(\Sigma,D)$ is independent of all
later runs, and $D<\infty$ almost surely. We put
\[
 \nu_\sgn(d)=\Prob(\Sigma=\sgn,\,D=d),\qquad Q_\sgn(N)=\Prob(\Sigma=\sgn,\,D\ge N),
 \qquad \pi_\sgn=\Prob(\Sigma=\sgn).
\]
Except in Conjecture~\ref{rem:aging}, the law of the initial run does not depend on
$N$. A fresh start of sign $\sgn_0$ has $\nu_{\sgn_0}=f_{\sgn_0}$ and
$\nu_{\osgn_0}=0$. The fair start is a fresh start of random sign, each sign with
probability $\frac12$.

Let $A_N$ be the number of plus steps among steps $1,\dots,N$, and
$P_N(k)=\Prob(A_N=k)$ for $0\le k\le N$. For a fresh start of sign $\sgn$ we write
$F^\sgn_N(k)$ for $P_N(k)$, with $F^\sgn_N(k)=0$ for $k\notin\{0,\dots,N\}$. We also
use
\[
 \Pi^\minus_m(j)=F^\minus_m(m-j),
\]
the probability that a fresh minus walk has exactly $j$ minus steps among its first
$m$ steps.

\subsection{Hypotheses}\label{sec:hypotheses}

The hypotheses carry fixed names, and every statement says which it uses.
\begin{itemize}[leftmargin=5em,labelsep=0.8em]
\item[(RV$_\alpha$)] $T_\plus$ is regularly varying with index $-\alpha$, where
$0<\alpha<1$, and $T_\minus(n)/T_\plus(n)\to b$ for some $b\in(0,\infty)$.
\item[(SPAN)] For each sign, the support of $X_\sgn$ is not contained in any
arithmetic progression $c+h\Z$ with $h\ge2$.
\item[(SRT$_\sgn$)] $u_\sgn(n)\sim\dfrac{\sin\pi\alpha}{\pi}\,\dfrac1{nT_\sgn(n)}$ as
$n\to\infty$, with $u_\sgn$ the renewal mass function of Section~\ref{sec:kernels}.
\item[(LB$_\sgn$)] There is $C$ with $f_\sgn(n)\le C\,T_\sgn(n)/n$ for all $n\ge1$.
\end{itemize}
Under (RV$_\alpha$) both tails are regularly varying with index $-\alpha$. Unequal
indices, or $b\in\{0,\infty\}$, give a degenerate limit of $A_N/N$ and are not
treated. (SPAN) is the aperiodicity needed for local limit theorems.
(SRT$_\sgn$) is the strong renewal theorem for the runs of sign $\sgn$. (LB$_\sgn$)
is Doney's local condition. It implies (SRT$_\sgn$) by Lemma~\ref{lem:cd}(iii)
below. A hypothesis written with the subscript $\pm$ is assumed for both signs.

\emph{Constants.} Unindexed $C$ and $c$ are positive constants that depend only on
the laws of the runs and may change from one use to the next. A constant written
$C_\delta$ also depends on $\delta$.

\subsection{Scales, stable densities and the Lamperti density}\label{sec:scales}

Under (RV$_\alpha$) choose an increasing function $a:(0,\infty)\to(0,\infty)$ with
$t\,T_\plus(a(t))\to1$ as $t\to\infty$, and put $a_m=a(m)$. So $a$ is regularly
varying with index $1/\alpha$. For $t>0$ let
\[
 m_N(t)=\lfloor t/T_\plus(N)\rfloor .
\]
Then $T_\plus(a_{m_N(t)})\sim T_\plus(N)/t$, and $a_{m_N(t)}/N\to t^{1/\alpha}$ uniformly for $t$
in compact subsets of $(0,\infty)$ (Lemma~\ref{lem:rv-facts}).

Put $\Gam=\Gamma(1-\alpha)$, and let $g$ be the density of the positive stable
variable $Z$ with $\E e^{-\lambda Z}=e^{-\Gam\lambda^\alpha}$. Its L\'evy measure is
$\alpha r^{-\alpha-1}dr$ on $(0,\infty)$, with tail $r^{-\alpha}$. For $t>0$ let
\[
 p_t(z)=t^{-1/\alpha}g\bigl(zt^{-1/\alpha}\bigr),
\]
the density at time $t$ of the stable subordinator with Laplace exponent
$t\Gam\lambda^\alpha$. The covering densities are
\begin{equation}\label{eq:covering}
 Q_t(x)=\int_0^xp_t(z)(x-z)^{-\alpha}\,dz,\qquad
 Q^b_t(y)=b\int_0^yp_{bt}(z)(y-z)^{-\alpha}\,dz .
\end{equation}
The density $g$ is continuous and bounded, and $g(w)\sim\alpha w^{-1-\alpha}$ as
$w\to\infty$~\cite[Section~XVII.6]{feller1971}. Hence
\begin{equation}\label{eq:p-bound}
 p_t(z)\le C\min\bigl\{t^{-1/\alpha},\,t\,z^{-1-\alpha}\bigr\}\qquad(t,z>0).
\end{equation}
For $0<x<1$ let
\begin{equation}\label{eq:Lambda}
 \Lambda(x)=\int_0^\infty\bigl[p_{bt}(1-x)\,Q_t(x)+p_t(x)\,Q^b_t(1-x)\bigr]\,dt .
\end{equation}
The Lamperti density $\La$ is defined in~\eqref{eq:lamperti-intro}. With
$t=\log((1-x)/x)$ and $\beta=\log b$ it reads
\begin{equation}\label{eq:L-cosh}
 \La(x)=\frac{\sin\pi\alpha}{\pi}\,\frac{1+\cosh t}{\cosh(\alpha t-\beta)+\cos\pi\alpha}.
\end{equation}
To check~\eqref{eq:L-cosh}, write $w=(1-x)/x=e^t$. Then
$x^{\alpha-1}(1-x)^{\alpha-1}=x^{2\alpha-2}w^{\alpha-1}$, the denominator
of~\eqref{eq:lamperti-intro} is $x^{2\alpha}bw^\alpha\cdot2[\cosh(\alpha t-\beta)+\cos\pi\alpha]$,
and $x^{-2}w^{-1}=(1+w)^2/w=2+2\cosh t$. We also use, for $0\le y<1$,
\begin{equation}\label{eq:Phi}
 \Phi_b(y)=\frac{b^2(1-y)^{\alpha-1}}{b^2(1-y)^{2\alpha}+2b\cos(\pi\alpha)\,(1-y)^\alpha y^\alpha+y^{2\alpha}},
\end{equation}
so $\Phi_b(0)=1$, $\Phi_b(y)=1+O(y^\alpha)$ as $y\to0$, $\Phi_b>0$, and
$\La(1-y)=\frac{\sin\pi\alpha}{\pi b}\,y^{\alpha-1}\Phi_b(y)$.

\subsection{Renewal kernels}\label{sec:kernels}

For each sign let $X^\sgn_1,X^\sgn_2,\dots$ be independent copies of $X_\sgn$,
$S^\sgn_m=X^\sgn_1+\dots+X^\sgn_m$ with $S^\sgn_0=0$, and $u^\sgn_m=f_\sgn^{*m}$ the
law of $S^\sgn_m$, so $u^\sgn_0=\delta_0$. The renewal mass function is
\[
 u_\sgn(j)=\sum_{m\ge0}u^\sgn_m(j),\qquad j\ge0,
\]
so $u_\sgn(0)=1$. For $m\ge1$ and $k\ge1$ let
\begin{equation}\label{eq:q-cover}
 q^\sgn_m(k)=\Prob\bigl(S^\sgn_{m-1}<k\le S^\sgn_m\bigr)=\sum_{r<k}u^\sgn_{m-1}(r)\,T_\sgn(k-r),
\end{equation}
the probability that the $m$-th run of sign $\sgn$ covers level $k$ in the time of
its sign. The second form holds because $X^\sgn_m$ is independent of $S^\sgn_{m-1}$.
For each $k\ge1$ exactly one $m$ has $S^\sgn_{m-1}<k\le S^\sgn_m$, so
$\sum_{m\ge1}q^\sgn_m(k)=1$. We write $\omega_\sgn(i)=\min\{1,1/(iT_\sgn(i))\}$ for
$i\ge1$ and $\omega_\sgn(0)=1$. Under (SRT$_\sgn$) we have $u_\sgn\le C\omega_\sgn$.

\subsection{Exact identities}\label{sec:identities}

The following identities hold for every $N$ and need no hypothesis. They are
proved in~\cite[Section~2]{pemmasaniF}, and we include the short proofs. The
computations of Section~\ref{sec:verify} check them against a sum over all sign
strings.

\begin{proposition}[the fresh-plus identity]\label{prop:fresh-identity}
For $k,j\ge1$ with $k+j=N$,
\begin{equation}\label{eq:fresh}
 F^\plus_N(k)=\sum_{m\ge1}u^\minus_m(j)\,q^\plus_{m+1}(k)+\sum_{m\ge1}u^\plus_m(k)\,q^\minus_m(j),
\end{equation}
and $F^\plus_N(N)=T_\plus(N)$, $F^\plus_N(0)=0$. Exchanging the signs, for $k,j\ge1$,
$F^\minus_N(k)=\sum_{m\ge1}[u^\plus_m(k)q^\minus_{m+1}(j)+u^\minus_m(j)q^\plus_m(k)]$,
with $F^\minus_N(0)=T_\minus(N)$ and $F^\minus_N(N)=0$.
\end{proposition}

\begin{proof}
For a fresh plus start the $i$-th plus run fills the steps in
$(S^\plus_{i-1}+S^\minus_{i-1},\,S^\plus_i+S^\minus_{i-1}]$ and the $i$-th minus run
fills those in $(S^\plus_i+S^\minus_{i-1},\,S^\plus_i+S^\minus_i]$. These intervals
partition $\{1,2,\dots\}$. If step $N$ lies in the $i$-th plus run, the first $N$
steps contain exactly $S^\minus_{i-1}$ minus steps. So $A_N=k$ means
$S^\minus_{i-1}=j$, and then step $N$ lies in that run exactly when
$S^\plus_{i-1}<k\le S^\plus_i$. The two events concern independent sequences, so
this case has probability $u^\minus_{i-1}(j)q^\plus_i(k)$. If step $N$ lies in the
$i$-th minus run, then $A_N=k$ means $S^\plus_i=k$, and step $N$ lies in that run
exactly when $S^\minus_{i-1}<j\le S^\minus_i$, with probability
$u^\plus_i(k)q^\minus_i(j)$. The cases are disjoint. For $j\ge1$ the plus case with
$i=1$ vanishes, and $m=i-1$ in the plus case and $m=i$ in the minus case
give~\eqref{eq:fresh}. For $k=N$ only the plus case with $i=1$ survives, with
probability $T_\plus(N)$, and $k=0$ is impossible. The fresh minus start is the same
with the signs exchanged.
\end{proof}

\begin{proposition}[the first-run decomposition]\label{prop:first-run}
For every initial run and every $0\le k\le N$,
\[
 P_N(k)=Q_\plus(N)\ones_{\{k=N\}}+Q_\minus(N)\ones_{\{k=0\}}
 +\sum_{d<N}\nu_\plus(d)\,F^\minus_{N-d}(k-d)+\sum_{d<N}\nu_\minus(d)\,F^\plus_{N-d}(k).
\]
In particular $P_N(N)=Q_\plus(N)$ and $P_N(0)=Q_\minus(N)$.
\end{proposition}

\begin{proof}
We condition on the initial run. If $\Sigma=\plus$ and $D\ge N$, then $A_N=N$. If
$\Sigma=\plus$ and $D=d<N$, steps $d+1,\dots,N$ are the first $N-d$ steps of a fresh
minus walk independent of $(\Sigma,D)$, so $A_N=d+A'$ with $A'$ of law
$F^\minus_{N-d}$. The case $\Sigma=\minus$ is the same. For $d<N$
Proposition~\ref{prop:fresh-identity} gives $F^\minus_{N-d}(N-d)=0$ and
$F^\plus_{N-d}(0)=0$, so only the atoms survive at $k=N$ and $k=0$.
\end{proof}

\begin{proposition}[the final-run identity]\label{prop:final-run}
For a fresh plus start, $N\ge1$ and $0\le j<N$,
\begin{equation}\label{eq:final-explicit}
\begin{split}
 P_N(N-j)={}&\sum_{r=0}^{j}u^\minus_r(j)\,\E\bigl[T_\plus(N-j-S^\plus_r);\,S^\plus_r\le N-j-1\bigr]\\
 &+\sum_{m=1}^{j}T_\minus(m)\sum_{r=0}^{j-m}u^\minus_r(j-m)\,u^\plus_{r+1}(N-j).
\end{split}
\end{equation}
\end{proposition}

\begin{proof}
We split according to the run that contains step $N$, the final run. The
$(r+1)$-th plus run starts at step $S^\plus_r+S^\minus_r+1$. It is the final run and
$A_N=N-j$ exactly when $S^\minus_r=j$ and $X^\plus_{r+1}\ge N-j-S^\plus_r\ge1$. This
gives the term $r$ of the first sum. The $(r+1)$-th minus run starts at step
$S^\plus_{r+1}+S^\minus_r+1$. If it is the final run and starts at step $N-m+1$, then
$A_N=N-j$ means $S^\plus_{r+1}=N-j$ and $S^\minus_r=j-m$, which needs $1\le m\le j$,
and $X^\minus_{r+1}\ge m$. This gives the second sum. Exactly one run is final, and
$u^\minus_r(i)=0$ for $r>i$.
\end{proof}

For the computations we use a third form of the law. For a fresh plus start let
$A(k,j)$ and $B(k,j)$ be the probabilities that a plus run, respectively a minus run,
starts right after exactly $k$ plus steps and $j$ minus steps. A plus run is preceded
by $m$ runs of each sign and a minus run by $m$ plus runs and $m-1$ minus runs, so
\[
 A(k,j)=\sum_{m\ge0}u^\plus_m(k)u^\minus_m(j),\qquad B(k,j)=\sum_{m\ge1}u^\plus_m(k)u^\minus_{m-1}(j).
\]
Splitting on the start of the final run gives, with $j=N-k$,
\begin{equation}\label{eq:run-sum}
 F^\plus_N(k)=\sum_{k'<k}A(k',j)\,T_\plus(k-k')+\sum_{j'<j}B(k,j')\,T_\minus(j-j').
\end{equation}
Every term is nonnegative.
